% Generated draft: requires mathematical and source review.
\documentclass[11pt]{article}
\usepackage[margin=25mm]{geometry}
\usepackage{fontspec}
\setmainfont{Latin Modern Roman}
\usepackage{amsmath,amssymb,mathtools,unicode-math}
\setmathfont{Latin Modern Math}
\usepackage{enumitem,needspace,xurl,hyperref}
\setlength{\parindent}{0pt}
\setlength{\parskip}{5pt}
\setlength{\emergencystretch}{3em}
\newcommand{\N}{\mathbb{N}}
\newcommand{\Z}{\mathbb{Z}}
\newcommand{\Q}{\mathbb{Q}}
\newcommand{\R}{\mathbb{R}}
\newcommand{\C}{\mathbb{C}}
\newcommand{\F}{\mathbb{F}}
\newcommand{\A}{\mathbb{A}}
\newcommand{\PP}{\mathbb{P}}
\newcommand{\E}{\mathbb{E}}
\newcommand{\eps}{\varepsilon}
\newcommand{\dd}{\,\mathrm{d}}
\begin{document}
\section*{UnsolvedMath ID: 30003603}
\textbf{Source problem code:} catalogue entry
\subsection{Definitions and mathematical statement}
Let \(M\cong\Z^d\) be a lattice, let \(M_{\R}=M\otimes_{\Z}\R\), and let \(P\subset M_{\R}\) be a full-dimensional lattice polytope. Its Ehrhart polynomial counts lattice points in dilates: \(L_P(n)=|nP\cap M|\) for integers \(n\geq0\). Define its Ehrhart \(h^*\)-polynomial and \(h^*\)-vector by
\[
\sum_{n\geq0}L_P(n)t^n=\frac{h_P^*(t)}{(1-t)^{d+1}},\qquad h_P^*(t)=\sum_{i=0}^{s}h_i^*t^i,\qquad (h_0^*,\ldots,h_s^*),
\]
where \(s=\deg h_P^*\leq d\). The coefficients are nonnegative integers and \(h_0^*=1\). The polynomial is \emph{real-rooted} if every complex zero is real; the constant polynomial satisfies this condition vacuously. Its coefficient vector is \emph{log-concave} if
\[
(h_i^*)^2\geq h_{i-1}^*h_{i+1}^*\qquad(1\leq i\leq s-1).
\]
It is \emph{unimodal} if there exists \(k\in\{0,\ldots,s\}\) such that \(h_0^*\leq\cdots\leq h_k^*\geq\cdots\geq h_s^*\).

For context, standard lattice-relative meanings of the geometric properties mentioned in the source are as follows. The polytope is \emph{smooth} if, at every vertex, the primitive lattice directions of its incident edges form a basis of \(M\). It is \emph{spanning} if
\[
\operatorname{span}_{\Z}\{x-y:x,y\in P\cap M\}=M.
\]
It is \emph{very ample} if, for every vertex \(v\),
\[
\left\{\sum_{x\in P\cap M}a_x(x-v):a_x\in\Z_{\geq0}\right\}
=\operatorname{cone}(P-v)\cap M,
\]
where \(\operatorname{cone}(P-v)\) denotes the nonnegative real cone generated by \(P-v\).

\textbf{Source question.} Can the implication diagram relating Ehrhart \(h^*\)-vector properties, including unimodality, to geometric properties such as smoothness, very ampleness, and spanning be extended to include real-rootedness of \(h_P^*(t)\) or log-concavity of its coefficients?

For a specified class \(\mathcal C\) of lattice polytopes, an arrow \(A\Longrightarrow B\) would mean
\[
\forall P\in\mathcal C,\qquad A(P)\Longrightarrow B(P).
\]
However, the supplied record does not reproduce the original diagram, enumerate all its predicates and arrows, or specify its full hypotheses and intended class \(\mathcal C\). The definitions above provide standard terminology for interpreting the question; they do not determine those missing choices. In particular, the intended ambient lattice convention for lower-dimensional polytopes is unspecified. Consequently, the source supplies a research question about extending an unspecified diagram, rather than a uniquely determined conjectural implication.
\subsection{Short English statement}
The question asks how real-rootedness and log-concavity of Ehrhart \(h^*\)-polynomials fit into an existing diagram connecting coefficient properties and geometric properties of lattice polytopes. The original diagram and its hypotheses are missing, so no particular new implication can be stated as the intended conjecture. The supplied open label requires independent verification.
\subsection{Sources}
\begin{itemize}
\item Supplied local UnsolvedMath registry record, permanent ID \(30003603\), title ``Real-Rootedness and Log-Concavity of Ehrhart Vectors,'' exact source problem code \texttt{catalogue entry}, in the Oberwolfach Reports Open Problems collection. No individual proposer is identified; the supplied proposed year is 2017.
\item Supplied report citation: \emph{Mini-Workshop: Lattice Polytopes: Methods, Advances, Applications} (2018), Oberwolfach Reports, pp. 2659--2701, DOI \texttt{10.4171/owr/2017/44}. Supplied source link: \url{https://doi.org/10.4171/owr/2017/44}. The record describes the reports as documentation of workshops and their open-problem discussions. The citation's 2018 date and the record's 2017 proposed year are retained without reconciliation.
\item Supplied literature source: Ferroni--Higashitani, ``examples and counterexamples,'' \url{https://arxiv.org/2307.10852}.
\item Supplied literature source: Braun, ``unimodality problems,'' \url{https://arxiv.org/1505.07377}.
\item Supplied literature source: ``original OWR report,'' \url{https://doi.org/10.4171/OWR/2017/44}.
\item Literature assessment dated \texttt{2026-08-21}: the supplied curation labels the problem open and reports isolated real-rooted and log-concave families, many counterexamples to naive implications, and no known general structural extension of the unimodality diagram to real-rooted or log-concave \(h^*\)-vectors. This is supplied dated curation derived from the OWR material, not an independently verified literature assessment. Its time-sensitive status and cited claims require independent verification.
\item The registry labels the statement \texttt{corrected\_verified} and records a manual semantic, typographic, and LaTeX review. That supplied editorial assessment does not resolve the missing diagram or constitute mathematical verification of this draft.
\end{itemize}
\end{document}
